% ======================================================================
% User Guide for the compass_sim Project — Version V78
% ======================================================================
\documentclass[11pt,a4paper]{article}

\usepackage[utf8]{inputenc}
\usepackage[T1]{fontenc}
\usepackage[english]{babel}
\usepackage{geometry}
\geometry{margin=2.5cm}
\usepackage{graphicx}
\usepackage{booktabs}
\usepackage{longtable}
\usepackage{amsmath}
\usepackage{amssymb}
\usepackage{siunitx}
\usepackage{hyperref}
\usepackage{xcolor}
\usepackage{listings}

\definecolor{codegray}{rgb}{0.95,0.95,0.95}
\definecolor{codeblue}{rgb}{0.1,0.2,0.6}
\definecolor{darkgray}{rgb}{0.3,0.3,0.3}

\lstset{
    backgroundcolor=\color{codegray},
    basicstyle=\small\ttfamily,
    breaklines=true,
    columns=fullflexible,
    commentstyle=\color{darkgray}\itshape,
    frame=single,
    keepspaces=true,
    keywordstyle=\color{codeblue}\bfseries,
    language=bash,
    showstringspaces=false
}

\hypersetup{
    colorlinks=true,
    linkcolor=codeblue,
    urlcolor=codeblue,
    pdftitle={User Guide: Compass-Needle Lattice Simulator (compass_sim) - Version V78},
    pdfauthor={J. P. Sinnecker}
}

\newcommand{\code}[1]{\texttt{#1}}

\newcount\hours
\newcount\minutes
\newcount\temp
\hours=\time
\divide\hours by 60
\minutes=\time
\temp=\hours
\multiply\temp by 60
\advance\minutes by -\temp
\def\compileTime{\ifnum\hours<10 0\fi\the\hours:\ifnum\minutes<10 0\fi\the\minutes}

\title{\textbf{User Guide: Compass-Needle Lattice Simulator (\code{compass\_sim})}\\[1ex]
\large Documentation and Reference for Version V78}
\author{\textbf{J.~P.~Sinnecker}}
\date{\today\ at \compileTime}

\begin{document}

\maketitle

\begin{abstract}
\noindent
This document is a comprehensive user guide and reference manual for the \code{compass\_sim} simulation engine (Version V78). The simulator models the two-dimensional collective dynamics of magnetic compass needles interacting via long-range dipolar fields. By resolving the full Newton's second law for rotation including needle inertia and viscous air damping, the simulator provides a realistic, physical platform for studying hysteresis loops, domain structures, and dynamic magnetic avalanches on square, triangular, and honeycomb lattices. This guide covers the physical model, configuration parameters, field sweep modes, demagnetization protocols, GPU acceleration options, and practical command-line examples.
\end{abstract}

\tableofcontents
\newpage

% ======================================================================
\section{Physics \& Theoretical Model}
\label{sec:physics}
% ======================================================================

The \code{compass\_sim} simulator resolves the real Newtonian rotational equations of motion for a two-dimensional lattice of classical magnetic dipoles (compass needles). For each needle $i$, the dynamics are governed by:

\begin{equation}
I \ddot{\theta}_i = \tau_i^{\text{dip}} + \tau_i^{\text{ext}} - b \dot{\theta}_i + \tau_i^{\text{noise}}
\label{eq:motion}
\end{equation}

\noindent
where the parameters are defined as follows:
\begin{itemize}
    \item $\theta_i$ is the in-plane angular orientation of the $i$-th needle relative to the $+x$ axis.
    \item $I$ is the needle's physical moment of inertia (\si{kg.m^2}).
    \item $\tau_i^{\text{dip}}$ is the mechanical torque exerted on needle $i$ by the total dipolar magnetic field $\vec{B}_i^{\text{dip}}$ generated by all other needles in the lattice.
    \item $\tau_i^{\text{ext}} = m (\hat{u}_i \times \vec{B}^{\text{ext}})_z$ is the torque exerted by the uniform external magnetic field $\vec{B}^{\text{ext}}$, where $m$ is the magnetic moment of the needle and $\hat{u}_i = (\cos\theta_i, \sin\theta_i)$ is its orientation vector.
    \item $b$ is the viscous air damping coefficient (\si{N.m.s/rad}), which governs dissipation.
    \item $\tau_i^{\text{noise}}$ is a random thermal noise torque, active only during simulated annealing runs.
\end{itemize}

The dipolar magnetic field $\vec{B}_i^{\text{dip}}$ generated at the center of needle $i$ by all other needles $j \neq i$ is computed using the full anisotropic dipolar formula:

\begin{equation}
\vec{B}_i^{\text{dip}} = \frac{\mu_0}{4\pi} \sum_{j \neq i} \frac{3(\vec{m}_j \cdot \hat{r}_{ij})\hat{r}_{ij} - \vec{m}_j}{r_{ij}^3}
\end{equation}

\noindent
where $\vec{r}_{ij} = \vec{r}_i - \vec{r}_j$, $r_{ij} = |\vec{r}_{ij}|$, and $\hat{r}_{ij} = \vec{r}_{ij}/r_{ij}$.

\subsection{Automatic Physical Parameter Calibration}
If not explicitly overridden in the CLI, the simulator automatically calibrates the physical properties of the needles by assuming each is a flat rectangular steel sheet of thickness $d$ (default $0.40\text{ mm}$), width $w = 0.22\,L$, and length $L = 2R \cdot \text{needle\_frac}$, where $R$ is the circumscribed circle radius of the lattice cell and $\text{needle\_frac}$ is the fractional needle size:
\begin{itemize}
    \item \textbf{Mass}: Computed using a standard carbon steel density $\rho \approx 7850\text{ kg/m}^3$:
    \begin{equation}
    \text{mass} = \rho \cdot L \cdot w \cdot d
    \end{equation}
    \item \textbf{Moment of Inertia ($I$)}: Assumes a uniform flat plate rotating about its center:
    \begin{equation}
    I = \frac{1}{12} \text{mass} \cdot (L^2 + w^2)
    \end{equation}
    \item \textbf{Magnetic Moment ($m$)}: Assumes the steel is magnetized to saturation ($B_{\text{sat}} = 2.0\text{ T}$, corresponding to $M_s = 1.59 \times 10^6\text{ A/m}$) along its longitudinal axis:
    \begin{equation}
    m = M_s \cdot \text{volume} = \frac{B_{\text{sat}}}{\mu_0} (L \cdot w \cdot d)
    \end{equation}
\end{itemize}

\subsection{Quality Factor and Dynamical Regime}
The degree of damping is characterized by the dimensionless quality factor $Q$:

\begin{equation}
Q = \frac{\omega_0 I}{b}
\end{equation}

\noindent
with $Q \gg 1$ underdamped (needles oscillate visibly before settling) and $Q \ll 1$ overdamped (monotonic relaxation). The transition between these regimes occurs at the critical value $Q_c = 0.5$. The default value of \code{--damping} ($5 \times 10^{-8}$ N$\cdot$m$\cdot$s/rad) places a standard tabletop needle at $Q \approx 25$, in the strongly underdamped regime.

\subsection{Integration Scheme and Time Step}
Equation~\eqref{eq:motion} is integrated numerically using a second-order \textbf{Velocity-Verlet} algorithm. To maintain numerical stability while ensuring physical correctness across diverse damping regimes, the integration time step $\Delta t$ is computed dynamically as a fraction of the needle's natural oscillation period $T_0$:

\begin{equation}
\Delta t = \text{dt\_factor} \cdot T_0 = \text{dt\_factor} \cdot \frac{2\pi}{\omega_0}
\end{equation}

\noindent
where $\text{dt\_factor}$ defaults to $0.05$, and $\omega_0 = \sqrt{m B_{\text{ref}}/I}$ is the natural frequency of a single needle oscillating in the reference dipolar field of its nearest neighbor:

\begin{equation}
B_{\text{ref}} = \frac{\mu_0}{4\pi}\frac{2m}{r_{\mathrm{nn}}^3}
\end{equation}

\noindent
with $r_{\mathrm{nn}} = 2R$ representing the center-to-center separation distance.

% ======================================================================
\section{Lattice Geometry Options}
\label{sec:geometries}
% ======================================================================

The spatial layout of the needles significantly shapes the local dipolar energy landscape. The simulator supports three distinct two-dimensional lattice geometry structures (configured via the \code{--geometry} flag):

\begin{enumerate}
    \item \textbf{Square Lattice} (\code{--geometry square}): A standard rectangular grid of $N$ rows and $M$ columns. Nearest neighbors are aligned along orthogonal axes at distance $r_{\mathrm{nn}} = 2R$. This layout has a coordination number $z=4$ (excluding boundaries).
    \item \textbf{Triangular Lattice} (\code{--geometry triangular}): An equilateral triangular lattice, packing needles more densely. Each internal needle has $z=6$ nearest neighbors located at $60^{\circ}$ angular intervals, which introduces geometric frustration.
    \item \textbf{Honeycomb Lattice} (\code{--geometry honeycomb}): A hexagonal grid resembling graphene. This structure features a lower packing density and coordination number $z=3$, resulting in different micro-domain dynamics.
\end{enumerate}

% ======================================================================
\section{External Field Modes}
\label{sec:modes}
% ======================================================================

The time-dependence of the external magnetic field $B_{\text{ext}}(t)$ is controlled via the \code{--field\_mode} argument. The supported modes are summarized in Table~\ref{tab:field_modes}.

\begin{table}[htbp]
\centering
\caption{Summary of available external field modes in \code{compass\_sim} V78.}
\label{tab:field_modes}
\vspace{1ex}
\begin{tabular}{lp{8.5cm}p{4.5cm}}
\toprule
\textbf{Mode} & \textbf{Description} & \textbf{Key Parameters} \\
\midrule
\code{static} & Applies a constant external field in a fixed direction. & \code{--B\_ext}, \code{--phi\_ext} \\
\code{hysteresis} & Sweeps the field through a full major hysteresis loop: $0 \to +B_{\max} \to 0 \to -B_{\max} \to 0 \to +B_{\max}$. & \code{--B\_ext} (sets $B_{\max}$), \code{--phi\_ext}, \code{--hyst\_spacing}, \code{--hyst\_log\_k} \\
\code{sine} & Applies a continuous sinusoidal field: $B_{\text{ext}}(t) = B_{\max}\sin(2\pi f t)$. & \code{--B\_ext} (sets $B_{\max}$), \code{--field\_freq} \\
\code{pulse} & Turns on field $B_{\max}$ after a delay for a set duration. & \code{--B\_ext}, \code{--field\_delay}, \code{--t\_pulse} \\
\code{step\_pos} & Instantly steps the field from $0$ to $+B_{\max}$ at $t = 0$. & \code{--B\_ext} \\
\code{step\_neg} & Instantly steps the field from $0$ to $-B_{\max}$ at $t = 0$. & \code{--B\_ext} \\
\code{forc} & Runs a sequence of First-Order Reversal Curves to map sub-loop magnetizations. & \code{--forc\_n\_curves}, \code{--forc\_Br\_min}, \code{--forc\_rate}, \code{--forc\_t\_ramp\_up} \\
\bottomrule
\end{tabular}
\end{table}

\subsection{First-Order Reversal Curves (FORC)}
The \code{forc} mode is a sophisticated characterization technique that sweeps the field as follows:
\begin{enumerate}
    \item Saturate the lattice in a strong positive field $+B_{\max}$.
    \item Ramp down to a reversal field $B_r$ (which varies systematically between curves).
    \item Sweep the field back up from $B_r$ to $+B_{\max}$, recording the magnetization along the way.
\end{enumerate}
This yields a family of minor curves that can be post-processed to calculate the joint distribution of local coercive and interaction fields. Version V78 supports constant sweep rate control (\code{--forc\_rate}), ensuring that the sweep rate $dB/dt$ remains identical across all curves regardless of their depth.

% ======================================================================
\section{Demagnetization Protocols}
\label{sec:demag}
% ======================================================================

To prepare the system in a clean, unbiased, isotropic state before starting the primary simulation, the lattice can be demagnetized using the \code{--demag} flag. The simulator implements two distinct protocols:

\subsection{Biaxial Rotational Demagnetization (\code{--demag rotational} or \code{--demag on})}
Instead of a standard uniaxial alternating field (which leaves a strong directional texture perpendicular to the shaking axis), this protocol rotates the external field vector continuously in the $xy$-plane while decaying its magnitude linearly to zero over a specified number of cycles:

\begin{align}
B_x(t) &= B_{\max} \left(1.0 - \frac{t}{t_{\text{demag}}}\right) \cos(2\pi f_{\text{demag}} t) \\
B_y(t) &= B_{\max} \left(1.0 - \frac{t}{t_{\text{demag}}}\right) \sin(2\pi f_{\text{demag}} t)
\end{align}

\noindent
where $t_{\text{demag}} = \text{demag\_cycles} / f_{\text{demag}}$. This circular dragging force erases directional memory in the lattice, resulting in an isotropic initial state.

\subsection{Simulated Thermal Annealing (\code{--demag anneal})}
This protocol decreases the magnitude of a linear shaking field while applying a decaying stochastic torque $\tau_i^{\text{noise}}$ directly to the needles, mimicking a high initial thermal temperature that cool down to zero. The random fluctuations allow the needle configurations to escape shallow metastable energy barriers and find the global dipolar ground state.

% ======================================================================
\section{System Requirements \& Acceleration}
\label{sec:acceleration}
% ======================================================================

Because calculating complete dipolar interactions has a computational complexity of $\mathcal{O}(K^2)$ for $K = N \times M$ needles, simulations on large lattices can become CPU-bound. The simulator dynamically adapts to your hardware:

\begin{itemize}
    \item \textbf{CPU Mode (Default)}: Powered by highly-vectorized NumPy operations. Ideal for grids up to $16 \times 16$.
    \item \textbf{GPU Mode (\code{--gpu 1})}: Accelerated by CuPy. This mode compiles custom CUDA kernels at runtime to perform needle-to-needle distance and field computations in parallel. Highly recommended for grid sizes of $30 \times 30$ (900 dipoles) or larger.
\end{itemize}

% ======================================================================
\section{Command Line Parameters Reference}
\label{sec:parameters}
% ======================================================================

The \code{compass\_sim} simulator provides a highly configurable interface. The parameters are categorized into six functional groups:

\subsection{Lattice Geometry Options}
These options configure the spatial dimensions and layout of the needle array:
\begin{description}
    \item[\code{--N <int>}] Number of rows in the lattice grid. Default: \code{8}.
    \item[\code{--M <int>}] Number of columns in the lattice grid. Default: \code{8}.
    \item[\code{--R <float>}] Radius of the circumscribed circle enclosing each needle (\si{m}). The center-to-center separation distance between nearest neighbor needles is $2R$. Default: \code{0.025} (\si{2.5\ cm}).
    \item[\code{--needle\_frac <float>}] Fractional needle length relative to cell diameter $2R$, bounded in $[0.0, 0.8]$. The physical needle length is $L = 2R \cdot \text{needle\_frac}$. Default: \code{0.80}. Values outside this range are clamped.
    \item[\code{--geometry \{square,triangular,honeycomb\}}] The physical layout structure of the grid. Default: \code{square}.
\end{description}

\subsection{Physics \& Material Properties}
These parameters dictate the mechanical and magnetic properties of individual needles, integration time step constraints, and boundary conditions:
\begin{description}
    \item[\code{--moment <float>}] Magnetic moment $m$ of each needle (\si{A.m^2}). If omitted, it is calculated from the steel volume ($L \times w \times d$) magnetized to saturation.
    \item[\code{--inertia <float>}] Physical moment of inertia $I$ (\si{kg.m^2}) of each needle. If omitted, it is computed assuming a uniform thin rectangular sheet rotating about its center.
    \item[\code{--pivot\_radius <float>}] Radius of the cylindrical pivot at the center of the needle (\si{m}). Default: \code{0.0} (no pivot).
    \item[\code{--pivot\_thickness <float>}] Thickness/height of the cylindrical pivot (\si{m}). Default: \code{0.0}.
    \item[\code{--pivot\_density <float>}] Density of the pivot material (\si{kg/m^3}). Default: \code{8500.0} (brass).
    \item[\code{--pivot\_mass <float>}] Explicit mass of the central pivot (\si{kg}). If supplied, it overrides density and dimension calculations.
    \item[\code{--needle\_thickness <float>}] Thickness $d$ of the needle sheet (\si{m}). Used for moment/inertia auto-calibration. Default: \code{0.0004} (\si{0.40\ mm}).
    \item[\code{--steel\_density <float>}] Density of the needle material (\si{kg/m^3}). Default: \code{7850.0} (carbon steel).
    \item[\code{--steel\_Bsat <float>}] Saturation magnetic flux density $B_{\text{sat}}$ (\si{T}). Default: \code{2.0}. Used for calculating magnetic moment ($M_s = B_{\text{sat}} / \mu_0$).
    \item[\code{--damping <float>}] Viscous damping coefficient $b$ (\si{N.m.s/rad}). Controls energy dissipation. Default: \code{5.00e-08}.
    \item[\code{--damping\_noise <float>}] Relative variation amplitude for per-needle damping. The damping of each needle $i$ is randomized as $b_i = \max(0, b \cdot (1 + \text{damping\_noise} \cdot U[-1, 1]))$. Default: \code{0.0}.
    \item[\code{--dt\_factor <float>}] Fraction of the natural oscillation period $T_0$ to use as the integration time step $\Delta t$. Lower values increase precision and stability. Default: \code{0.05}.
    \item[\code{--noise <float>}] Angular noise amplitude (\si{rad}) applied to initial needle orientations. Default: \code{1.5}.
    \item[\code{--seed <int>}] Seed for the pseudorandom number generator. If omitted, the current time in nanoseconds is used.
    \item[\code{--pbc \{0,1\}}] Toggle periodic boundary conditions (\code{0} = disabled, \code{1} = enabled). Default: \code{0}.
    \item[\code{--pbc\_images <int>}] Number of image shells used to compute long-range periodic dipolar terms. Default: \code{1}.
    \item[\code{--gpu \{0,1\}}] Enable GPU acceleration via CuPy. Highly recommended for lattice sizes exceeding $30 \times 30$. Default: \code{0}.
    \item[\code{--progress\_bar \{0,1\}}] Toggle the visual terminal progress bar. Default: \code{1}.
    \item[\code{--halo\_mode \{order,domains\}}] Coloring scheme for needle visualization. \code{order} shades based on individual deviation, while \code{domains} colors by localized orientation cluster. Default: \code{order}.
    \item[\code{--domain\_tol <float>}] Tolerance angle in degrees used to group needles into discrete domains in \code{domains} halo mode. Default: \code{15.0}.
\end{description}

\subsection{External Magnetic Field Options}
These parameters define the field components and profiles:
\begin{description}
    \item[\code{--field\_mode \{static,hysteresis,sine,pulse,step\_pos,step\_neg,forc\}}] The time-varying protocol for the applied external field $\vec{B}^{\text{ext}}(t)$. Default: \code{static}.
    \item[\code{--B\_ext <float>}] Magnitude of the external magnetic field (\si{T}). Default: \code{0.0}.
    \item[\code{--phi\_ext <float>}] Angle of the external magnetic field (\si{degrees}) in the $xy$-plane, relative to the $+x$ axis. Default: \code{0.0}.
    \item[\code{--ext\_Bx <float>}] Explicit $x$-component of the external magnetic field (\si{T}). Overrides \code{--B\_ext} and \code{--phi\_ext}.
    \item[\code{--ext\_By <float>}] Explicit $y$-component of the external magnetic field (\si{T}). Overrides \code{--B\_ext} and \code{--phi\_ext}.
    \item[\code{--field\_delay <float>}] Time delay (\si{s}) before the active field excitation protocol starts. Default: \code{0.0}.
    \item[\code{--field\_freq <float>}] Frequency $f$ (\si{Hz}) for the sinusoidal external field. Default: \code{1.0}.
    \item[\code{--t\_pulse <float>}] Duration of the field pulse (\si{s}). Only relevant when \code{--field\_mode} is \code{pulse}.
    \item[\code{--torque\_tol <float>}] Relative torque threshold used to determine if the lattice has reached mechanical equilibrium (early-stop trigger). Default: \code{1e-3}.
\end{description}

\subsection{FORC \& Demagnetization Parameters}
Parameters specific to advanced field sweeps and pre-simulation preparation:
\begin{description}
    \item[\code{--forc\_n\_curves <int>}] Number of minor reversal curves to simulate in FORC mode. Default: \code{30}.
    \item[\code{--forc\_Br\_min <float>}] Minimum reversal field $B_{r,\min}$ (\si{T}). Defaults to $-B_{\max}$ if unspecified.
    \item[\code{--forc\_t\_sat <float>}] Duration (\si{s}) the system is held in positive saturation $+B_{\max}$ prior to starting each curve. Default: \code{0.05} (used if \code{--forc\_rate} is not specified).
    \item[\code{--forc\_t\_ramp\_down <float>}] Time (\si{s}) to ramp from positive saturation down to the reversal field $B_r$. Default: \code{0.10} (used if \code{--forc\_rate} is not specified).
    \item[\code{--forc\_t\_ramp\_up <float>}] Time (\si{s}) to sweep back from $B_r$ up to $+B_{\max}$. Default: \code{0.20} (used if \code{--forc\_rate} is not specified).
    \item[\code{--forc\_sweep \{both,up,down\}}] Trace increasing field curves, decreasing field curves, or both. Default: \code{both}.
    \item[\code{--forc\_rate <float>}] Target constant sweep rate $R$ (\si{T/s}). If specified, the sweep durations are scaled dynamically to enforce a constant rate $dB/dt = R$ across all loops. Default: \code{None}.
    \item[\code{--demag \{off,on,rotational,anneal\}}] Selects a pre-relaxation protocol to demagnetize the lattice: \code{off} (none), \code{on}/\code{rotational} (rotational decaying field), or \code{anneal} (simulated thermal annealing). Default: \code{off}.
    \item[\code{--demag\_freq <float>}] Rotation/oscillation frequency (\si{Hz}) of the demagnetization field. Default: \code{2.0}.
    \item[\code{--demag\_cycles <int>}] Number of full decay cycles during demagnetization. Default: \code{20}.
    \item[\code{--demag\_temp <float>}] Noise multiplier factor representing initial thermal energy during simulated annealing. Default: \code{1.0}.
    \item[\code{--demag\_delay <float>}] Settle time (\si{s}) between the end of demagnetization and the start of the primary simulation. Default: \code{0.0}.
    \item[\code{--hyst\_spacing \{linear,log\}}] Spacing distribution style of field steps in hysteresis sweeps. Default: \code{linear}.
    \item[\code{--hyst\_log\_k <float>}] Concentration parameter controlling spacing density for logarithmic hysteresis sweeps. Default: \code{5.0}.
\end{description}

\subsection{Execution \& Stop Controls}
\begin{description}
    \item[\code{--t\_sim <float>}] Total physical time (\si{s}) allocated for the simulation run. Default: \code{2.0}.
    \item[\code{--t\_sim\_full \{0,1\}}] If set to \code{1}, forces the simulator to run for the full duration of \code{--t\_sim}, bypassing all early physical stopping triggers (e.g., full alignment $S=1.0$ or mechanical rest). Useful for maintaining uniform time-series datasets. Default: \code{0}.
\end{description}

\subsection{Visualization \& Output Options}
These controls manage visual rendering, frame rates, and data export:
\begin{description}
    \item[\code{--video <path>}] Filename or absolute path to export the simulation video (compiled in MP4 format).
    \item[\code{--frame\_every <int>}] Saves one video frame every $N$ integration steps. Default: \code{5}.
    \item[\code{--make\_images \{0,1\}}] Enable or disable generating summary PNG plots at the end of the simulation. Default: \code{1}.
    \item[\code{--fps <int>}] Frame rate (frames per second) of the exported MP4 video. Default: \code{24}.
    \item[\code{--dpi <int>}] Dot-per-inch resolution of saved frames. Default: \code{120}.
    \item[\code{--keep\_frames}] Keep the temporary directory of PNG frames after compilation. Default: \code{False} (frames deleted).
    \item[\code{--hyst\_autoscale \{0,1\}}] Enable or disable autoscale on the axes of the hysteresis plots. Default: \code{1}.
    \item[\code{--csv\_order \{t,B\}}] Order of the export columns in the time-series CSV file: \code{t} writes time first, \code{B} writes field first. Default: \code{t}.
\end{description}

% ======================================================================
\section{Output Files \& Data Formats}
\label{sec:outputs}
% ======================================================================

The simulator generates standard files to facilitate analysis and visualization:

\subsection{Visual Plots (PNG)}
When \code{--make\_images 1} is active, the following summary plots are generated in the working directory:
\begin{itemize}
    \item \code{compass\_initial.png}: Shows the initial angular orientation configuration of the needles.
    \item \code{compass\_equilibrium.png}: Displays the final, relaxed configuration of the needle lattice.
    \item \code{compass\_comparison.png}: Displays a side-by-side comparison of the initial and final states of the needle array.
    \item \code{compass\_order\_param.png}: Plots the temporal evolution of the lattice-wide order parameter $S(t) = \frac{1}{K} \left| \sum_{j=1}^{K} e^{i\theta_j} \right|$.
    \item \code{hysteresis\_loop.png} or \code{forc\_curves.png}: Generated during field sweeps. Plots the projected lattice magnetization $M_{\text{proj}}$ as a function of the external field $B_{\text{ext}}$.
    \item \code{sine\_field.png}: Time-series of both the sinusoidal external field and the lattice's projected magnetization.
\end{itemize}

\subsection{Data Logs (CSV)}
The simulator logs numerical data to CSV format. The filename is determined by the field mode and options:
\begin{itemize}
    \item \textbf{General Log (\code{compass\_field\_log.csv} or \code{<video\_name>.csv})}:
    By default, column structure is:
    \begin{lstlisting}[backgroundcolor=\color{white},frame=none]
t_s, B_applied_T, M_proj, S
    \end{lstlisting}
    If \code{--csv\_order B} is set, the structure changes to:
    \begin{lstlisting}[backgroundcolor=\color{white},frame=none]
B_applied_T, M_proj, S, t_s
    \end{lstlisting}
    
    \item \textbf{Hysteresis Loop (\code{hysteresis\_loop.csv})}:
    Logs the sweep trajectory. Columns: \code{t\_s}, \code{B\_T}, \code{M\_proj}, \code{S}.
    
    \item \textbf{FORC Sweeps (\code{forc\_curve.csv})}:
    Logs families of minor loops. Includes the reversal field $B_r$:
    \begin{lstlisting}[backgroundcolor=\color{white},frame=none]
t_s, B_T, M_proj, S, B_r
    \end{lstlisting}
    
    \item \textbf{Sinusoidal Sweeps (\code{sine\_field.csv})}:
    Logs the sinusoidal response. Columns: \code{t\_s}, \code{B\_T}, \code{M\_proj}, \code{S}.
\end{itemize}

\subsection{Simulation Videos (MP4)}
If \code{--video <name>} is specified, a high-quality video is rendered showing the dynamic spatial evolution of the lattice. Needle colors dynamically correspond to their instantaneous angle or local domain affinity depending on the coloring mode.


\newpage
% ======================================================================
\section{Step-by-Step Practical Examples}
\label{sec:examples}
% ======================================================================

Here are standard recipes for running common simulation procedures.

\subsection{Example 1: Basic Equilibrium Relaxation (No Field)}
Relax a square lattice from a randomized state to analyze domain formation:
\begin{lstlisting}
python3 compass.py --t_sim 2.0 --noise 2.5
\end{lstlisting}

\subsection{Example 2: Major Hysteresis Loop}
Simulate a full hysteresis cycle for a $10 \times 10$ triangular lattice with a maximum field of $1.5\text{ mT}$ applied at an angle of $30^{\circ}$:
\begin{lstlisting}
python3 compass.py --geometry triangular --N 10 --M 10 --field_mode hysteresis --B_ext 1.5e-3 --phi_ext 30 --t_sim 5.0 --damping 5e-7
\end{lstlisting}

\subsection{Example 3: Rotational Demagnetization \& Pulse}
First, rotational-demagnetize the lattice at $5\text{ Hz}$ over $40$ cycles, let the needles settle for a $0.2\text{ s}$ delay, and then apply a $100\ \mu\text{T}$ pulse:
\begin{lstlisting}
python3 compass.py --field_mode pulse --B_ext 100e-6 --phi_ext 90 --demag rotational --demag_freq 5.0 --demag_cycles 40 --demag_delay 0.2 --t_pulse 0.5 --t_sim 1.5 --video pulse_demag.mp4
\end{lstlisting}

\subsection{Example 4: High-Performance GPU Run}
Simulate a large $30 \times 30$ lattice (900 needles) with GPU acceleration to create a smooth high-resolution video:
\begin{lstlisting}
python3 compass.py --N 30 --M 30 --gpu 1 --t_sim 2.5 --video large_grid.mp4
\end{lstlisting}

% ======================================================================
\section{Running Sweep Campaigns}
\label{sec:sweeps}
% ======================================================================

For scientific studies requiring parametric sweeps across damping coefficients, lattice structures, and random seeds, use the orchestration campaign script:

\begin{lstlisting}
python3 sweep_damping_hysteresis.py --out_dir ./damping_campaign --gpu 1
\end{lstlisting}

\noindent
This script schedules a 120-run grid campaign (8 damping/Q values $\times$ 3 geometries $\times$ 5 seeds) and includes full resume functionality (\code{--resume}) which skips already completed CSV logs.

\newpage
% ======================================================================
\section{Summary Reference Table of Parameters}
\label{sec:summary_table}
% ======================================================================

This section provides a quick lookup for all command-line interface arguments available in \code{compass\_sim}.

\begin{longtable}{lp{3.5cm}lp{6.0cm}}
\caption{Complete CLI parameter guide for \code{compass\_sim} Version V78.}\\
\toprule
\textbf{Parameter} & \textbf{Type / Options} & \textbf{Default} & \textbf{Brief Description} \\
\midrule
\endfirsthead
\multicolumn{4}{c}%
{{\bfseries Table \thetable\ (continued from previous page)}} \\
\toprule
\textbf{Parameter} & \textbf{Type / Options} & \textbf{Default} & \textbf{Brief Description} \\
\midrule
\endhead
\bottomrule
\multicolumn{4}{r}{{Continued on next page...}} \\
\endfoot
\bottomrule
\endlastfoot

\multicolumn{4}{l}{\textit{Lattice Geometry Options}} \\
\midrule
\code{--N} & \code{int} & \code{8} & Number of rows in the lattice grid. \\
\code{--M} & \code{int} & \code{8} & Number of columns in the lattice grid. \\
\code{--R} & \code{float} & \code{0.025} & Cell circumscribed radius (\si{m}); center spacing is $2R$. \\
\code{--needle\_frac} & \code{float} & \code{0.80} & Needle size ratio relative to $2R$, range $[0.0, 0.8]$. \\
\code{--geometry} & \code{square, triangular, honeycomb} & \code{square} & Physical spatial layout arrangement. \\
\midrule

\multicolumn{4}{l}{\textit{Physics \& Material Properties}} \\
\midrule
\code{--moment} & \code{float} & \code{None} & Needle magnetic moment $m$ (\si{A.m^2}). Calculated from geometry if omitted. \\
\code{--inertia} & \code{float} & \code{None} & Needle moment of inertia $I$ (\si{kg.m^2}). Calculated from geometry if omitted. \\
\code{--pivot\_radius} & \code{float} & \code{0.0} & Radius of the cylindrical pivot (\si{m}). \\
\code{--pivot\_thickness} & \code{float} & \code{0.0} & Thickness/height of the central cylindrical pivot (\si{m}). \\
\code{--pivot\_density} & \code{float} & \code{8500.0} & Density of pivot material (\si{kg/m^3}). \\
\code{--pivot\_mass} & \code{float} & \code{None} & Explicit mass of the central pivot (\si{kg}). Overrides density/dimensions. \\
\code{--needle\_thickness} & \code{float} & \code{0.0004} & Thickness $d$ (\si{m}) used for moment/inertia auto-calibration. \\
\code{--steel\_density} & \code{float} & \code{7850.0} & Density of needle material (\si{kg/m^3}). \\
\code{--steel\_Bsat} & \code{float} & \code{None} & Saturation flux density $B_{\text{sat}}$ (\si{T}). Auto-calibrates $M_s$ if omitted (default: $2.0$ T). \\
\code{--damping} & \code{float} & \code{5.00e-08} & Viscous damping coefficient $b$ (\si{N.m.s/rad}). \\
\code{--damping\_noise} & \code{float} & \code{0.0} & Relative variation factor in damping coefficient per needle. \\
\code{--dt\_factor} & \code{float} & \code{0.05} & Integration step factor $\Delta t / T_0$ for Velocity-Verlet. \\
\code{--noise} & \code{float} & \code{1.5} & Angular noise amplitude (\si{rad}) for initial orientations. \\
\code{--seed} & \code{int} & \code{None} & Pseudorandom number generator seed. Uses time-ns if omitted. \\
\code{--pbc} & \code{0, 1} & \code{0} & Toggle periodic boundary conditions. \\
\code{--pbc\_images} & \code{int} & \code{1} & Periodic replication shell summation depth. \\
\code{--gpu} & \code{0, 1} & \code{0} & Toggle GPU acceleration via CuPy. \\
\code{--progress\_bar} & \code{0, 1} & \code{1} & Toggle visual interactive CLI progress bar. \\
\code{--halo\_mode} & \code{order, domains} & \code{order} & Halo coloring mode for needle rendering. \\
\code{--domain\_tol} & \code{float} & \code{15.0} & Angle threshold (degrees) for cluster-grouping domains. \\
\midrule

\multicolumn{4}{l}{\textit{External Magnetic Field Options}} \\
\midrule
\code{--field\_mode} & \code{static, hysteresis, sine, pulse, step\_pos, step\_neg, forc} & \code{static} & The excitation profile for the applied field $\vec{B}^{\text{ext}}(t)$. \\
\code{--B\_ext} & \code{float} & \code{0.0} & External field magnitude limit (\si{T}). \\
\code{--phi\_ext} & \code{float} & \code{0.0} & Applied field angle (degrees) relative to $+x$. \\
\code{--ext\_Bx} & \code{float} & \code{None} & Direct $x$-component of the external magnetic field (\si{T}). \\
\code{--ext\_By} & \code{float} & \code{None} & Direct $y$-component of the external magnetic field (\si{T}). \\
\code{--field\_delay} & \code{float} & \code{0.0} & Delay duration (\si{s}) before applying field protocol. \\
\code{--field\_freq} & \code{float} & \code{1.0} & Sinusoidal AC field frequency (\si{Hz}). \\
\code{--t\_pulse} & \code{float} & \code{None} & Pulse duration (\si{s}) under pulse mode. \\
\code{--torque\_tol} & \code{float} & \code{1e-3} & Rotational equilibrium early stop tolerance. \\
\midrule

\multicolumn{4}{l}{\textit{FORC \& Demagnetization Parameters}} \\
\midrule
\code{--forc\_n\_curves} & \code{int} & \code{30} & Number of First-Order Reversal Curves to trace. \\
\code{--forc\_Br\_min} & \code{float} & \code{None} & Minimum reversal field limit (\si{T}). Defaults to $-B_{\max}$. \\
\code{--forc\_t\_sat} & \code{float} & \code{None} & Hold time at $+B_{\max}$ before each curve (\si{s}); internal default $0.05$~s. \\
\code{--forc\_t\_ramp\_down} & \code{float} & \code{None} & Ramp time from saturation to $B_r$ (\si{s}); internal default $0.10$~s. \\
\code{--forc\_t\_ramp\_up} & \code{float} & \code{None} & Sweep time back from $B_r$ to $+B_{\max}$ (\si{s}); internal default $0.20$~s. \\
\code{--forc\_sweep} & \code{both, up, down} & \code{both} & Curve sweep directions to save and record in CSV. \\
\code{--forc\_rate} & \code{float} & \code{None} & Constant $dB/dt$ rate (\si{T/s}) for dynamic curve duration scaling. \\
\code{--demag} & \code{off, on, rotational, anneal} & \code{off} & Selected demagnetization method prior to simulation. \\
\code{--demag\_freq} & \code{float} & \code{2.0} & Frequency (\si{Hz}) of demagnetization fields. \\
\code{--demag\_cycles} & \code{int} & \code{20} & Number of field decay cycles for demagnetization. \\
\code{--demag\_temp} & \code{float} & \code{1.0} & Stochastic noise factor for simulated annealing. \\
\code{--demag\_delay} & \code{float} & \code{0.0} & Rest settle time (\si{s}) after demagnetization ends. \\
\code{--hyst\_spacing} & \code{linear, log} & \code{linear} & Spacing distribution type for hysteresis field steps. \\
\code{--hyst\_log\_k} & \code{float} & \code{5.0} & Log-spacing density parameter for hysteresis. \\
\midrule

\multicolumn{4}{l}{\textit{Execution \& Stop Controls}} \\
\midrule
\code{--t\_sim} & \code{float} & \code{2.0} & Total simulation physical time (\si{s}). \\
\code{--t\_sim\_full} & \code{0, 1} & \code{0} & Enforce full $t_{\text{sim}}$ duration execution without early stops. \\
\midrule

\multicolumn{4}{l}{\textit{Visualization \& Output Options}} \\
\midrule
\code{--video} & \code{str} & \code{None} & Filepath to save rendered MP4 animation. \\
\code{--frame\_every} & \code{int} & \code{5} & Save frame image every $N$ integration steps. \\
\code{--make\_images} & \code{0, 1} & \code{1} & Enable summary and comparison PNG plot exports. \\
\code{--fps} & \code{int} & \code{24} & Frames per second of compiled video. \\
\code{--dpi} & \code{int} & \code{120} & Dot-per-inch resolution of frames. \\
\code{--keep\_frames} & flag & \code{False} & Keep raw PNG frames folder after video compile. \\
\code{--hyst\_autoscale} & \code{0, 1} & \code{1} & Enable autoscale on the axes of hysteresis plots. \\
\code{--csv\_order} & \code{t, B} & \code{t} & Column structure in field log CSV file (t-first or B-first). \\

\end{longtable}

\end{document}
